% ===========================================================================
%  Exercises - paper implementation (look "paper")
% ---------------------------------------------------------------------------
%  Plain amsthm environments, an italic run-in "Answer." for worked solutions
%  and a plain a), b), ... list for subtasks. No boxes, no colour.
%
%  Selected by style/exercises.tex, which documents the commands.
% ===========================================================================

\theoremstyle{definition}
\newtheorem{paperexercise}{Exercise}[section]
\newtheorem{paperhomework}[paperexercise]{Homework}

% Compatibility aliases for the environment names of the boxed look
\newenvironment{myexercise}[2]{\begin{paperexercise}[#1]}{\end{paperexercise}}
\newenvironment{myhomework}[2]{\begin{paperhomework}[#1]}{\end{paperhomework}}

\NewDocumentCommand{\exercise}{m+m}{%
    \IfBlankTF{#1}
        {\begin{paperexercise}#2\end{paperexercise}}
        {\begin{paperexercise}[#1]#2\end{paperexercise}}%
}

\NewDocumentCommand{\exerciser}{mm+m}{%
    \IfBlankTF{#1}
        {\begin{paperexercise}\label{exc:#2}#3\end{paperexercise}}
        {\begin{paperexercise}[#1]\label{exc:#2}#3\end{paperexercise}}%
}

\NewDocumentCommand{\hw}{m+m}{%
    \IfBlankTF{#1}
        {\begin{paperhomework}#2\end{paperhomework}}
        {\begin{paperhomework}[#1]#2\end{paperhomework}}%
}

\NewDocumentCommand{\hwr}{mm+m}{%
    \IfBlankTF{#1}
        {\begin{paperhomework}\label{hw:#2}#3\end{paperhomework}}
        {\begin{paperhomework}[#1]\label{hw:#2}#3\end{paperhomework}}%
}

% ---------------------------------------------------------------------------
%  Answer blocks
% ---------------------------------------------------------------------------
\newenvironment{exercisepf}{%
    \par\addvspace{0.4em}%
    \noindent\textit{Answer.}\ \ignorespaces
}{%
    \par\addvspace{0.6em}%
}

\newenvironment{exercisepartanswer}{%
    \par\addvspace{0.2em}%
    \noindent\textit{Answer.}\ \ignorespaces
}{%
    \par\addvspace{0.3em}%
}

\newlist{exerciseparts}{enumerate}{1}
\setlist[exerciseparts]{
    label=\textup{(\alph*)},
    leftmargin=2.2em,
    itemsep=0.8em,
    topsep=0.6em,
    parsep=0pt
}

\NewDocumentEnvironment{ecparts}{}{
    \begin{exerciseparts}
}{
    \end{exerciseparts}
}

\NewDocumentCommand{\ecpart}{m+m}{%
    \item #1

    \begin{exercisepartanswer}
        #2
    \end{exercisepartanswer}
}

\NewDocumentCommand{\ec}{m+m+m}{%
    \exercise{#1}{#2}
    \begin{exercisepf}
        #3
    \end{exercisepf}
}

\NewDocumentCommand{\ecs}{m+m+m}{%
    \exercise{#1}{#2}
    \begin{ecparts}
        #3
    \end{ecparts}
}
